\section{Discussion}
\label{sec:discussion}

The model is best interpreted as a theory of second-best local green industrial policy rather than a complete theory of climate policy. A national or supranational carbon price can address the global emissions externality while local governments retain responsibility for place-based infrastructure, industrial development, and implementation. Our question is how those remaining local instruments interact when firms compete across jurisdictions.

The distinction between $s_i$ and $h_i$ matters economically. A targeted subsidy directly changes the private return to one firm's green investment. Shared infrastructure changes local production conditions and may also reduce emissions directly. The switching result therefore concerns the composition of a response, not simply the total amount of green support. At high rivalry a local government can find infrastructure attractive precisely because it supports competitiveness without entering a one-for-one subsidy contest.

The result also clarifies why the two firm investment margins should not be collapsed into a single generic investment variable. Conventional investment is not introduced to create a mechanical resource trade-off; investment costs are separable. Instead, it provides an endogenous adjustment margin that amplifies policy shocks through the product market. The no-$x$ benchmark shows that removing this margin changes the qualitative strategic response.

Several limits should be kept in view. Fixed plant locations exclude the discrete bidding and relocation forces emphasized in the location-competition literature. The territorial emissions objective is a reduced-form representation of delegated climate targets and does not capture global carbon damages. Infrastructure is summarized by one scalar rather than modeled as a network or congestible public input. Finally, the model is symmetric in its canonical form. These choices keep the mechanism transparent. Extensions are valuable only if they alter the switching logic or generate empirically distinguishable predictions; otherwise they add complexity without strengthening the contribution.

The empirical predictions in Section~\ref{sec:institutional} provide a way to distinguish the mechanism from generic policy imitation. The key prediction is an interaction: the response of local shared infrastructure to rival targeted support should become more positive as product-market rivalry increases, especially in sectors where shared infrastructure and green investment have meaningful productivity effects. Evidence only that neighboring jurisdictions copy one another's spending would not be sufficient.
